\documentclass[12pt]{article}
\usepackage[margin=1in]{geometry}
\usepackage{graphicx}
\usepackage{booktabs}
\usepackage{array}
\usepackage{longtable}
\usepackage{hyperref}
\usepackage{caption}
\usepackage{float}
\usepackage{titlesec}
\usepackage{enumitem}
\usepackage{parskip}
\usepackage{fancyhdr}
\usepackage{xcolor}

\hypersetup{
    colorlinks=true,
    linkcolor=black,
    urlcolor=blue
}

\pagestyle{fancy}
\fancyhf{}
\rhead{CS 4632 -- Milestone 3}
\lhead{Munir Gargour}
\rfoot{\thepage}

\graphicspath{{output/}}

\title{
    \textbf{CS 4632: Modeling and Simulation}\\
    \large Milestone 3: Complete Implementation and Testing\\[0.5em]
    \normalsize Seasonal Predator-Prey Simulation
}
\author{Munir Gargour}
\date{}

\begin{document}

\maketitle
\thispagestyle{fancy}
\tableofcontents
\newpage

% =============================================================================
\section{Implementation Summary}
% =============================================================================

\subsection{Overview of Current Implementation}

The Seasonal Predator-Prey Simulation is a continuous 2D agent-based model
implemented in Python. It simulates the population dynamics of rabbits (prey),
foxes (predators), and grass (resource) on a toroidal $500 \times 500$ spatial
plane across a full simulated year divided into four seasons (Spring, Summer,
Fall, Winter), each 250 simulation steps in length. The simulation is fully
parameterized via JSON configuration files and outputs time-series CSVs,
summary JSONs, and population dynamic charts for each run.

The core software components are:

\begin{itemize}
    \item \textbf{Animal (base class):} Manages position, energy, speed,
    metabolism, and the toroidal movement model.

    \item \textbf{Rabbit (subclass of Animal):} Implements Algorithm 1
    (Hunger-Driven Foraging) and Algorithm 3 (Reproduction Mode). Rabbits
    search for the nearest grass patch within a foraging radius using KD-Tree
    spatial queries, consume grass when within range, and enter reproduction
    mode when energy is sufficient.

    \item \textbf{Fox (subclass of Animal):} Implements Algorithm 2 (Predator
    Pursuit) and Algorithm 3 (Reproduction Mode). Foxes detect nearby rabbits
    within a vision radius, pursue and attempt to catch the nearest rabbit, and
    reproduce when energy exceeds a threshold.

    \item \textbf{Environment:} Manages all entities, executes per-step grass
    regeneration scaled by seasonal multipliers, and provides KD-Tree backed
    spatial query methods for neighbor lookups.

    \item \textbf{Simulation:} Orchestrates the main time-step loop, advances
    seasons every 250 steps, collects time-series and summary data, and writes
    output files.

    \item \textbf{run\_all.py:} Iterates over all JSON configs in the
    \texttt{configs/} directory, runs each simulation, and writes a master
    index JSON summarizing all runs.
\end{itemize}

\subsection{Features Completed in M3}

All features proposed in M2 have been implemented. The following were
completed for this milestone:

\begin{itemize}
    \item Four-season dynamics with per-season metabolism multipliers and grass
    growth multipliers applied each step.

    \item Toroidal boundary wrapping for continuous uninterrupted movement
    across the $500 \times 500$ plane.

    \item Long-phase multi-year tracking: Run 006 executes 2000 steps (two
    full simulated years) to observe multi-cycle behavior.

    \item Comprehensive automated data collection: per-step time-series CSV,
    per-run summary JSON, per-run config JSON copy, and population dynamic PNG
    charts.

    \item A master index JSON aggregating results from all 10 runs with timing
    and outcome metadata.

    \item 10 distinct simulation runs covering a variety of parameter
    configurations.
\end{itemize}

\subsection{Scope Changes from M2 Proposal}

Several parameter adjustments were made relative to the M2 baseline during M3
tuning to achieve ecologically plausible dynamics:

\begin{itemize}
    \item \textbf{Rabbit \texttt{max\_energy} reduced from 100 to 75.} With
    \texttt{max\_energy} at 100, rabbits frequently accumulated excess energy
    that could not be used productively. By setting the foraging threshold to
    \texttt{max\_energy} minus \texttt{energy\_per\_grass} (i.e., 65), no grass
    energy is wasted due to overflow. This made rabbit energy dynamics more
    responsive to available food and improved cycle regularity.

    \item \textbf{Fox \texttt{base\_metabolism} reduced from 1.0 to 0.8 and
    reproduction probability from 0.35 to 0.25.} These adjustments extended
    individual fox survival during periods of low rabbit density while
    preventing fox populations from growing large enough to collapse the rabbit
    population.

    \item \textbf{Fox \texttt{vision\_range} reduced from 30 to 14.} At range
    30 in a $500 \times 500$ space, foxes detected prey too easily at high
    rabbit densities, causing rapid rabbit collapse. Range 14 produces more
    realistic boom-bust-recovery cycles.

    \item \textbf{Base grass spawn rate raised from 5 to 8 patches per step.}
    This provides a more stable grass floor through Fall and Winter and sustains
    rabbit populations into the next Spring cycle.
\end{itemize}

\subsection{Architecture and Design}

The simulation uses a discrete time-step model with a fixed step size. Each
step proceeds as follows:

\begin{enumerate}
    \item Grass regeneration: $n_{new} = \text{round}(\text{base\_spawn\_rate}
    \times \text{grass\_mult})$ new patches are placed at random positions.
    \item All living rabbits execute their \texttt{step()} method (forage or
    reproduce).
    \item All living foxes execute their \texttt{step()} method (hunt or
    reproduce).
    \item Dead entities are removed from the entity lists.
    \item Current state is recorded to the time-series buffer.
    \item Season is advanced every 250 steps.
\end{enumerate}

Spatial queries are handled by \texttt{scipy.spatial.KDTree}, rebuilt each
step as needed for grass lookup and rebuilt on demand for animal neighbor
searches. The toroidal wrapping is implemented via modulo arithmetic on both
axes. All output files are written to the \texttt{output/} directory using the
naming convention \texttt{run\_XXX\_<type>.<ext>} where \texttt{XXX} is the
zero-padded run ID.

% =============================================================================
\section{Execution Documentation}
% =============================================================================

\subsection{Run Summary Table}

\begin{table}[H]
\centering
\caption{Summary of all 10 simulation runs}
\small
\begin{tabular}{clrrrrrcc}
\toprule
\textbf{Run} & \textbf{Description} & \textbf{Steps} & \textbf{Time (s)} &
\textbf{R\textsubscript{final}} & \textbf{F\textsubscript{final}} &
\textbf{R ext.} & \textbf{F ext.} \\
\midrule
001 & Baseline, all default parameters       & 1000 & 23.2 & 87  & 0  & No & Yes \\
002 & High rabbit count (200 rabbits)        & 1000 & 23.2 & 86  & 0  & No & Yes \\
003 & High fox count (20 foxes)              & 1000 & 16.4 & 81  & 0  & No & Yes \\
004 & Low initial grass (50 patches)         & 1000 & 21.4 & 90  & 0  & No & Yes \\
005 & High grass abundance (500 patches)     & 1000 & 38.2 & 133 & 0  & No & Yes \\
006 & Long run, 2 full years (2000 steps)    & 2000 & 52.2 & 82  & 0  & No & Yes \\
007 & Harsh winter (metabolism $\times$2.0)  & 1000 & 20.2 & 45  & 0  & No & Yes \\
008 & Mild seasons, uniform year-round       & 1000 & 18.7 & 161 & 10 & No & No  \\
009 & High entity count (300R, 15F, 300G)    & 1000 & 28.8 & 91  & 0  & No & Yes \\
010 & Different seed (seed=99), baseline     & 1000 & 20.4 & 89  & 0  & No & Yes \\
\bottomrule
\end{tabular}
\end{table}

Total wall time for all 10 runs: 262.7 seconds. All 10 runs completed
successfully. No rabbit extinction was observed in any run. Fox extinction
occurred in 9 out of 10 runs. Run 008 (mild seasons) is the only scenario
where foxes persisted to the end of the simulation ($F = 10$).

\subsection{Parameter Configurations}

Key parameters varied across runs (values differing from baseline are noted):

\begin{description}[leftmargin=2em]
    \item[Run 001 -- Baseline]
    \texttt{num\_rabbits}: 100, \texttt{num\_foxes}: 5,
    \texttt{initial\_grass}: 200, \texttt{spawn\_rate}: 8.
    Seasons: Spring (met=0.8, grs=2.0), Summer (met=1.0, grs=1.5),
    Fall (met=1.1, grs=0.7), Winter (met=1.4, grs=0.4).

    \item[Run 002 -- High Rabbit Count]
    \texttt{num\_rabbits}: \textbf{200}, all other parameters same as baseline.

    \item[Run 003 -- High Fox Count]
    \texttt{num\_foxes}: \textbf{20}, all other parameters same as baseline.

    \item[Run 004 -- Low Initial Grass]
    \texttt{initial\_grass}: \textbf{50}, all other parameters same as baseline.

    \item[Run 005 -- High Grass Abundance]
    \texttt{initial\_grass}: \textbf{500}, \texttt{spawn\_rate}: \textbf{13},
    all other parameters same as baseline.

    \item[Run 006 -- Long Run (2 Years)]
    \texttt{steps}: \textbf{2000}, all other parameters same as baseline.

    \item[Run 007 -- Harsh Winter]
    Winter season changed to: met=\textbf{2.0}, grs=\textbf{0.20}.
    All other parameters same as baseline.

    \item[Run 008 -- Mild Seasons]
    All seasonal multipliers changed to near-uniform values:
    Spring (met=0.95, grs=1.1), Summer (met=1.0, grs=1.0),
    Fall (met=1.05, grs=0.95), Winter (met=1.1, grs=0.9).

    \item[Run 009 -- High Entity Count]
    \texttt{num\_rabbits}: \textbf{300}, \texttt{num\_foxes}: \textbf{15},
    \texttt{initial\_grass}: \textbf{300}. Seasons same as baseline.

    \item[Run 010 -- Different Random Seed]
    \texttt{seed}: \textbf{99}, all other parameters same as baseline.
\end{description}

Parameters varied across the 10 runs include: initial rabbit count, initial
fox count, initial grass count, grass spawn rate, simulation duration,
seasonal metabolism multipliers, seasonal grass multipliers, and random seed,
covering well more than the required minimum of 3 distinct parameters.

\subsection{Execution Environment}

\begin{itemize}
    \item \textbf{Operating System:} Windows 11 Pro (10.0.26200)
    \item \textbf{Python Version:} 3.13
    \item \textbf{Key Libraries:} NumPy, SciPy (KDTree), Matplotlib
    \item \textbf{Performance:} Runs completed in 16 to 52 seconds each
    depending on entity count and step count. Total suite time approximately
    4 minutes 23 seconds.
\end{itemize}

\subsection{Issues Encountered}

During M3 execution the primary challenge was balancing predator and prey
parameters to produce ecologically plausible behavior. The key issues and
resolutions were:

\begin{itemize}
    \item \textbf{Fox overhunting:} With high vision range (30), foxes detected
    prey too easily at peak rabbit densities and drove rabbits to extinction
    within a single Spring. Resolution: vision range reduced to 14, reducing
    the expected number of rabbits detectable per step by approximately 78\%.

    \item \textbf{Fox extinction during lean periods:} Foxes died off before
    rabbit populations could recover because they lacked sufficient energy
    reserves between successful hunts. Resolution: fox \texttt{base\_metabolism}
    reduced from 1.0 to 0.8, extending survival without eating by 25\%. Fox
    reproduction probability was simultaneously reduced from 0.35 to 0.25 to
    prevent the lower metabolism from causing unchecked breeding.

    \item \textbf{Grass depletion leading to dampened cycles:} At peak rabbit
    populations, consumption greatly exceeded regeneration, causing grass to
    crash near zero. Rabbits then entered a survival-only state with
    insufficient energy to breed, flattening population cycles. Resolution:
    \texttt{base\_grass\_spawn\_rate} raised from 5 to 8 and rabbit
    \texttt{max\_energy} reduced from 100 to 75 to improve energy utilization
    efficiency per grass patch consumed.
\end{itemize}

% =============================================================================
\section{Data Collection Overview}
% =============================================================================

\subsection{Metrics Collected}

The following metrics are recorded automatically during each simulation run:

\textbf{Time-Series Data} (\texttt{run\_XXX\_timeseries.csv}, one row per step):
\begin{itemize}
    \item \texttt{step}: Simulation step number
    \item \texttt{year}: Simulated year number
    \item \texttt{season}: Current season name
    \item \texttt{rabbits}, \texttt{foxes}, \texttt{grass}: Living entity counts
    \item \texttt{rabbits\_born}, \texttt{foxes\_born}: Births this step
    \item \texttt{rabbits\_eaten}, \texttt{grass\_eaten}: Consumption events
    \item \texttt{metabolism\_mult}, \texttt{grass\_mult}: Active seasonal multipliers
\end{itemize}

\textbf{Summary Statistics} (\texttt{run\_XXX\_summary.json}):
\begin{itemize}
    \item Total steps, wall time, final and average entity counts
    \item Maximum and minimum observed rabbit and fox counts
    \item Total births, predation events, and grass consumed
    \item Boolean extinction flags for rabbits and foxes
\end{itemize}

\textbf{Configuration Snapshot} (\texttt{run\_XXX\_config.json}): Full copy of
the input configuration file archived with each run.

\textbf{Master Index} (\texttt{master\_index.json}): Aggregated outcome table
for all 10 runs including run ID, description, steps ran, wall time, final
counts, and output file references.

\subsection{Data Sample}

The following is an excerpt from \texttt{run\_001\_timeseries.csv} (first 4
steps):

\begin{table}[H]
\centering
\caption{Excerpt from run\_001\_timeseries.csv}
\small
\begin{tabular}{rrlrrrrrrrcc}
\toprule
\textbf{step} & \textbf{yr} & \textbf{season} & \textbf{R} & \textbf{F} &
\textbf{G} & \textbf{R born} & \textbf{F born} & \textbf{R eaten} &
\textbf{G eaten} & \textbf{met} & \textbf{grs} \\
\midrule
1 & 1 & Spring & 154 & 5 & 201 & 55 & 0 & 1 & 15  & 0.8 & 2.0 \\
2 & 1 & Spring & 177 & 5 & 190 & 78 & 0 & 1 & 42  & 0.8 & 2.0 \\
3 & 1 & Spring & 188 & 5 & 169 & 89 & 0 & 1 & 79  & 0.8 & 2.0 \\
4 & 1 & Spring & 192 & 5 & 153 & 93 & 0 & 1 & 111 & 0.8 & 2.0 \\
\bottomrule
\end{tabular}
\end{table}

The rapid rabbit growth in the first few Spring steps (100 to 192 in 4 steps)
reflects the low fox population and abundant grass creating ideal breeding
conditions. Grass is consumed far faster than it is regenerated in these early
steps due to the large and growing rabbit population.

\subsection{Visualizations}

Each run produces a two-panel population dynamics chart. The top panel shows
rabbit (blue) and fox (red) population counts over time with season bands
shaded in the background. The bottom panel shows grass patch counts over the
same period. All 10 charts are presented in Section~\ref{sec:results}.

% =============================================================================
\section{Preliminary Results}
\label{sec:results}
% =============================================================================

\subsection{Basic Statistics Across All Runs}

\begin{table}[H]
\centering
\caption{Aggregate statistics for all 10 runs}
\small
\begin{tabular}{lrrrrrrr}
\toprule
\textbf{Run} & \textbf{Avg R} & \textbf{Max R} & \textbf{Avg F} &
\textbf{Max F} & \textbf{Grass Eaten} & \textbf{R Born} & \textbf{F Born} \\
\midrule
001 & 188.6 & 544   & 20.4 & 54  & 9,429  & 2,273 & 54  \\
002 & 203.3 & 591   & 21.4 & 61  & 9,421  & 2,232 & 70  \\
003 & 109.4 & 289   & 31.6 & 73  & 9,430  & 2,506 & 67  \\
004 & 185.5 & 509   & 20.0 & 45  & 9,270  & 2,198 & 48  \\
005 & 302.1 & 883   & 33.4 & 104 & 15,486 & 3,962 & 134 \\
006 & 274.9 & 791   & 10.2 & 54  & 18,681 & 3,753 & 54  \\
007 & 172.0 & 544   & 18.5 & 54  & 8,652  & 2,059 & 53  \\
008 & 171.9 & 291   & 14.1 & 20  & 8,167  & 1,732 & 21  \\
009 & 250.4 & 1,014 & 19.9 & 79  & 9,530  & 2,350 & 70  \\
010 & 172.4 & 484   & 22.0 & 55  & 9,431  & 2,323 & 56  \\
\bottomrule
\end{tabular}
\end{table}

\subsection{Observed Patterns and Trends}

\subsubsection{Predator-Prey Lag Effect}

In all runs, fox population peaks consistently lag behind rabbit population
peaks by approximately 50 to 100 steps. This is the classic Lotka-Volterra
delay: rabbits must first accumulate in large numbers before foxes have
sufficient prey encounters to breed and grow their own population. This lag
is clearly visible in Figure~\ref{fig:run001}.

\begin{figure}[H]
    \centering
    \includegraphics[width=\textwidth]{run_001_population.png}
    \caption{Run 001 -- Baseline population dynamics. Fox peak lags rabbit
    peak by approximately 75 steps.}
    \label{fig:run001}
\end{figure}

\subsubsection{Seasonal Boom-Bust Cycles}

The baseline run (001) demonstrates a clear Spring boom where rabbits peak at
544 before foxes grow large enough (peak 54) to drive a crash. Rabbit
populations partially recover in subsequent seasons before declining again
through Winter as grass becomes scarce and metabolism costs rise. This
multi-cycle damped oscillation is visible across the majority of runs.

\subsubsection{Effect of Initial Population Size}

Comparing Run 001 (100 rabbits) to Run 002 (200 rabbits) shows that a higher
initial rabbit count accelerates fox population growth (avg foxes 20.4 vs
21.4, max foxes 54 vs 61) but produces similar long-term outcomes. The
additional initial rabbits are consumed in the first Spring wave and the
systems converge to comparable steady-state dynamics by Summer.

\begin{figure}[H]
    \centering
    \includegraphics[width=\textwidth]{run_002_population.png}
    \caption{Run 002 -- High initial rabbit count (200). Peak rabbit count
    reaches 591 vs 544 in the baseline.}
    \label{fig:run002}
\end{figure}

Run 003 (20 initial foxes) constrains rabbits much more sharply. Average rabbit
count drops to 109 compared to 188 in the baseline, and the maximum rabbit
count only reaches 289 vs 544. Higher initial fox pressure suppresses the early
Spring rabbit boom almost entirely.

\begin{figure}[H]
    \centering
    \includegraphics[width=\textwidth]{run_003_population.png}
    \caption{Run 003 -- High initial fox count (20). Rabbit population is
    heavily suppressed throughout, peaking at only 289.}
    \label{fig:run003}
\end{figure}

\subsubsection{Effect of Grass Availability}

Run 004 (50 initial grass) vs Run 005 (500 initial grass, spawn rate 13) shows
the strongest environmental contrast. With scarce initial grass, rabbits grow
more slowly and peak at only 509 compared to 883 in the high-grass run. Average
grass in Run 004 was 74 patches vs 259 patches in Run 005. The fox population
in Run 005 reached a maximum of 104, the highest of any run, driven by the
large rabbit population that abundant food resources enabled.

\begin{figure}[H]
    \centering
    \includegraphics[width=\textwidth]{run_004_population.png}
    \caption{Run 004 -- Low initial grass (50 patches). Slower rabbit growth
    and lower peak population due to early food scarcity.}
    \label{fig:run004}
\end{figure}

\begin{figure}[H]
    \centering
    \includegraphics[width=\textwidth]{run_005_population.png}
    \caption{Run 005 -- High grass abundance (500 initial, spawn rate 13).
    Rabbit peak of 883 and fox peak of 104 are the highest of any run.}
    \label{fig:run005}
\end{figure}

\subsubsection{Multi-Year Dynamics}

Run 006 ran for 2000 steps, covering two full simulated years. The second year
showed noticeably lower peak rabbit counts as the fox population established
from Year 1 suppressed the Year 2 Spring boom more aggressively. Average rabbit
count over 2000 steps was 274.9 but average fox count dropped to 10.2 due to
the extended period of fox absence after Year 1 extinction. This demonstrates
that the simulation does not recover a fox population once it collapses, as
expected for a closed system with no immigration.

\begin{figure}[H]
    \centering
    \includegraphics[width=\textwidth]{run_006_population.png}
    \caption{Run 006 -- Two full simulated years (2000 steps). Second-year
    dynamics show suppressed rabbit peaks and absent fox recovery.}
    \label{fig:run006}
\end{figure}

\subsubsection{Harsh Winter Effects}

Under doubled Winter metabolism (2.0$\times$) and severely reduced grass growth
(0.20$\times$), rabbit populations still survived but finished at only 45
individuals, the lowest final rabbit count of any run. The reduced grass
availability through Winter kept rabbit energy in a survival-only state,
suppressing reproduction and producing the most severe population decline of all
standard runs.

\begin{figure}[H]
    \centering
    \includegraphics[width=\textwidth]{run_007_population.png}
    \caption{Run 007 -- Harsh winter conditions. Rabbit count reaches a
    minimum of 17 during Winter, the most severe decline of any run.}
    \label{fig:run007}
\end{figure}

\subsubsection{Stable Coexistence Under Mild Conditions}

Run 008 is the only run where foxes survived to the end of the simulation
($F = 10$). With nearly uniform seasonal multipliers (metabolism ranging only
from 0.95 to 1.10 and grass from 0.90 to 1.10), the rabbit population never
crashed hard enough to starve the fox population out. Average rabbit count was
171.9 and fox count 14.1, both maintained without major oscillations. This
confirms that seasonal variation is the primary driver of the boom-bust cycles
and eventual fox extinction seen in other runs.

\begin{figure}[H]
    \centering
    \includegraphics[width=\textwidth]{run_008_population.png}
    \caption{Run 008 -- Mild uniform seasons. The only run with fox survival
    at end of simulation ($F = 10$), demonstrating stable coexistence.}
    \label{fig:run008}
\end{figure}

\subsubsection{High Entity Count}

Starting with 300 rabbits, 15 foxes, and 300 grass produced the highest single
rabbit peak of any run at 1,014 individuals. Despite this, the long-term
outcome was similar to lower-density runs with a final rabbit count of 91. The
larger initial populations amplified cycle amplitude but did not change the
qualitative pattern.

\begin{figure}[H]
    \centering
    \includegraphics[width=\textwidth]{run_009_population.png}
    \caption{Run 009 -- High entity count (300R, 15F, 300G). Rabbit population
    exceeds 1,000 individuals at peak, the highest of any run.}
    \label{fig:run009}
\end{figure}

\subsubsection{Seed Sensitivity}

Run 010 uses the same parameters as Run 001 but with seed 99 instead of 42.
Final rabbit count (89 vs 87) and average rabbits (172 vs 188) are very
similar, confirming that the qualitative dynamics are robust to random seed
variation. Differences in exact peak timing and amplitude are present but the
overall trajectory is consistent, which supports the validity of the model.

\begin{figure}[H]
    \centering
    \includegraphics[width=\textwidth]{run_010_population.png}
    \caption{Run 010 -- Alternative random seed (seed=99). Dynamics closely
    mirror Run 001, confirming model robustness to stochastic variation.}
    \label{fig:run010}
\end{figure}

\subsection{Performance Observations}

Simulation performance scales primarily with entity count and step count. The
baseline 1000-step run completes in approximately 23 seconds. Run 005, which
sustains very high rabbit populations (avg 302) throughout, took 38 seconds due
to the increased number of entity updates and spatial queries per step. Run 006
(2000 steps) took 52 seconds, confirming that step count is also a significant
contributor to wall time.

All runs completed without error or memory issues. The KD-Tree spatial query
approach implemented in M2 continues to scale well even at the maximum entity
count seen in Run 009 (peak of 1,014 rabbits plus up to 79 foxes simultaneously
active).

\subsection{Summary of Key Findings}

\begin{itemize}
    \item Predator-prey oscillation with seasonal lag is consistently observed
    across all 10 runs, validating the core model behavior.

    \item Rabbit populations are resilient: no rabbit extinction was observed in
    any of the 10 runs including the harsh winter scenario.

    \item Fox persistence depends strongly on seasonal regularity. In 9 of 10
    runs foxes eventually go extinct due to insufficient prey density during
    Winter in this closed-system environment.

    \item Grass availability is the primary indirect driver of cycle amplitude.
    Higher grass leads to larger rabbit peaks, larger fox peaks, and more
    dramatic crashes.

    \item The simulation demonstrates ecologically plausible Lotka-Volterra
    dynamics with the added complexity of seasonal forcing, making it a useful
    tool for exploring how environmental variability shapes predator-prey
    coexistence.
\end{itemize}

% =============================================================================
\section*{Submission Checklist}
% =============================================================================

\begin{itemize}
    \item[$\checkmark$] Completed at least 10 simulation runs (10 runs documented)
    \item[$\checkmark$] Documented all parameter variations (Section 2.2)
    \item[$\checkmark$] Included data samples in report (Section 3.2)
    \item[$\checkmark$] Created run summary table (Sections 2.1 and 4.1)
    \item[$\checkmark$] Documented scope changes from M2 (Section 1.3)
    \item[$\checkmark$] Included preliminary observations (Section 4.2)
    \item[$\checkmark$] Report is well-organized and clear
\end{itemize}

\end{document}
